\documentclass[11pt]{article}

\usepackage[margin=1in]{geometry}
\usepackage{amsmath, amssymb}
\usepackage{xcolor}
\usepackage{listings}
\usepackage{hyperref}
\usepackage{array}
\usepackage{booktabs}

\hypersetup{
    colorlinks=true,
    linkcolor=blue!60!black,
    urlcolor=blue!60!black,
    citecolor=blue!60!black
}

\lstdefinestyle{bugpython}{
  language=Python,
  basicstyle=\ttfamily\small,
  keywordstyle=\color{blue!60!black},
  stringstyle=\color{green!40!black},
  commentstyle=\color{gray!70!black},
  showstringspaces=false,
  breaklines=true,
  frame=single,
  rulecolor=\color{gray!30},
  columns=fullflexible,
  keepspaces=true
}

\title{SymPy Bug Report}
\author{}
\date{}

\begin{document}

\maketitle

\section{Bug 22: solveset returns 1 for acos(x) = 2*pi outside the principal acos range}

\subsection{Minimal Reproducer}

\medskip

\begin{lstlisting}[style=bugpython]
from sympy import Eq, S, N, acos, pi, solveset, symbols

x = symbols("x")
eq = Eq(acos(x), 2*pi)
sol = solveset(eq, x, domain=S.Complexes)

print("solution =", sol)
for s in sol:
    print("residual at", s, "=", N(acos(s) - 2*pi, 80))
\end{lstlisting}

\subsection{Observed Behavior}

\medskip

\begin{lstlisting}[style=bugpython]
solution = {1}
residual at 1 = -6.2831853071795864769252867665590057683943387987502116419498891846156328125724180
\end{lstlisting}

SymPy returns the candidate \(x = 1\).  Substituting that value into the
original equation gives
\[
  \operatorname{acos}(1) - 2\pi = -2\pi,
\]
so the returned value does not satisfy the equation.

\subsection{Expected Behavior}

The correct solution set over \(S.\mathrm{Complexes}\) is \(\varnothing\).  The
principal value of \(\operatorname{acos}(1)\) is \(0\), not \(2\pi\), and the principal branch
of \(\operatorname{acos}\) does not attain the real value \(2\pi\).

\subsection{Mathematical Explanation}

The equation being solved is
\[
  \operatorname{acos}(x) = 2\pi.
\]
It is tempting to apply \(\cos\) to both sides and infer
\[
  x = \cos(2\pi) = 1.
\]
That inference is invalid for the principal inverse cosine.  The identity
\(\cos(\operatorname{acos}(x)) = x\) holds on the appropriate domain, but \(\operatorname{acos}(\cos(y))\)
does not equal \(y\) for arbitrary \(y\).  Instead, \(\operatorname{acos}\) returns a
principal value.  In particular,
\[
  \operatorname{acos}(1) = 0,
\]
so \(x = 1\) satisfies \(\operatorname{acos}(x) = 0\), not \(\operatorname{acos}(x) = 2\pi\).  The reported
answer is therefore an extraneous solution introduced by treating a
principal-branch inverse as though it were a globally valid inverse.

\subsection{Source-Level Diagnosis}

\begin{sloppypar}
The diagnosis identifies the relevant solver code in
\nolinkurl{sympy/solvers/solveset.py}, specifically
\texttt{\_solveset} and \texttt{\_invert\_complex}.  The traced call path starts
from the complex-domain call to solve \(\operatorname{acos}(x) = 2\pi\).  The
equation is normalized to \texttt{acos(x) - 2*pi}.  Because this is not
classified as a \texttt{TrigonometricFunction} equation, \texttt{\_solveset}
takes the generic inversion path and calls the inverter on the normalized
equation.  The complex inverter removes the additive \(-2\pi\), obtaining
\texttt{acos(x) = 2*pi}, and then uses the generic inverse-function branch.
\end{sloppypar}

The key rule in \(\texttt{\_invert\_complex}\) is:

\begin{lstlisting}[style=bugpython]
if hasattr(f, 'inverse') and f.inverse() is not None and \
   not isinstance(f, TrigonometricFunction) and \
   not isinstance(f, HyperbolicFunction) and \
   not isinstance(f, exp):
    ...
    return _invert_complex(f.args[0],
                           imageset(Lambda(n, f.inverse()(n)), g_ys), symbol)
\end{lstlisting}

For \(f = \operatorname{acos}(x)\), the inverse method returns \(\cos\).  That
inverse is defined in \nolinkurl{sympy/functions/elementary/trigonometric.py}.
With target set \(\{2\pi\}\), this branch constructs an image set equivalent to
\(\{\cos(2\pi)\}\), which evaluates to \(\{1\}\).  The diagnosis further notes
that \texttt{domain\_check} only checks that the original expression is defined
at \(x = 1\); it does not verify that the residual is zero.  Thus the false
finite candidate survives even though
\(\operatorname{acos}(1) - 2\pi = -2\pi\).

A plausible fix direction supported by the diagnosis is to avoid using this
generic inverse branch unconditionally for principal inverse functions.  The
solver could handle inverse trigonometric functions with explicit
principal-range conditions, or it could filter finite inverted candidates by
checking that they satisfy the original equation.

\subsection{Additional Failing Instances}

The same mechanism appears for other real targets outside the principal range
of \(\operatorname{acos}\).  Applying \(\cos\) to those targets produces \(1\) or \(-1\), but
the principal values \(\operatorname{acos}(1) = 0\) and \(\operatorname{acos}(-1) = \pi\) do not equal the
requested targets.  The independent verification records the following
parametrized cases.

\begin{center}
\small
\renewcommand{\arraystretch}{1.3}
\begin{tabular}{@{}>{\raggedright\arraybackslash}p{0.44\linewidth}>{\raggedright\arraybackslash}p{0.23\linewidth}>{\raggedright\arraybackslash}p{0.23\linewidth}@{}}
\toprule
\textbf{Input / Expression} & \textbf{SymPy behavior} & \textbf{Expected behavior} \\
\midrule
\(\operatorname{acos}(x) = 2\pi\) & returns a set containing \(1\) & \(\varnothing\) \\
\(\operatorname{acos}(x) = -\pi\) & returns a set containing \(-1\) & \(\varnothing\) \\
\(\operatorname{acos}(x) = 3\pi\) & returns a set containing \(-1\) & \(\varnothing\) \\
\(\operatorname{acos}(x) = 4\pi\) & returns a set containing \(1\) & \(\varnothing\) \\
\(\operatorname{acos}(x) = -2\pi\) & returns a set containing \(1\) & \(\varnothing\) \\
\(\operatorname{acos}(x) = 5\pi\) & returns a set containing \(-1\) & \(\varnothing\) \\
\bottomrule
\end{tabular}
\end{center}

The expected behavior is the same across the family because each target is a
real value outside the principal real range of \(\operatorname{acos}\), while the returned
candidate is only a solution of the periodic cosine equation obtained after an
invalid global inversion step.

\subsection{Related Issues Assessment}

The deduplication evidence found related background on principal branches of
inverse trigonometric functions and the SymPy \texttt{solveset} contract, but no
public report of this exact reproducer or output.  The verdict is therefore a
new instance of a known failure mode: the broad issue of respecting principal
branches is known, but this exact \texttt{acos} case with the false result
\(\{1\}\) appears previously unreported and remains individually actionable as
a focused issue and regression test.

\subsection{Proposed SymPy Regression Test}

\medskip

\begin{lstlisting}[style=bugpython]
import pytest
from sympy import Eq, S, acos, pi, solveset, symbols, simplify

x = symbols("x")

CASES = [
    (2*pi, 1),
    (-pi, -1),
    (3*pi, -1),
    (4*pi, 1),
    (-2*pi, 1),
    (5*pi, -1),
]


@pytest.mark.parametrize("target,bad", CASES)
def test_solveset_acos_rejects_values_outside_principal_range(target, bad):
    sol = solveset(Eq(acos(x), target), x, domain=S.Complexes)

    assert sol == S.EmptySet
    assert simplify(acos(bad) - target) != 0
\end{lstlisting}

\end{document}
